\chapter{Python 爬虫基础：网页结构、HTTP、requests 与正则表达式}
\lstdefinestyle{pyD06}{language=Python,basicstyle=\ttfamily\footnotesize,aboveskip=3pt,belowskip=3pt}
\lstdefinestyle{txtD06}{basicstyle=\ttfamily\footnotesize,aboveskip=3pt,belowskip=3pt}
\noindent\key{本章速览}\quad 爬虫五步：\key{拼网址} → \key{取源码}（\texttt{requests.get(url, headers=headers).text}）→ \key{正则提取}（“文本A\texttt{(.*?)}文本B”+\texttt{re.S}）→ \key{清洗} → \key{输出存盘}；Selenium、反爬、PDF 解析见第 7 章。

% ======================================================================
\section{金融研究中的文本数据：四类来源与研究用途}

\subsection{为什么需要爬虫：从“硬数据”到“软信息”}
\key{硬数据}（股价、财务科目）之外，文本数据（年报措辞、新闻数量与语气、政策强度）度量\key{软信息}（soft information）：风险、不确定性、战略取向、投资者情绪。文本散落在网站上，需\term{网络爬虫}{web crawler / web scraper}自动访问网页、下载内容、提取字段。

\subsection{四类文本数据：公司文本、媒体舆情、平台商业、政策监管}
\begin{center}\footnotesize\renewcommand{\arraystretch}{1.0}
\begin{tabularx}{\linewidth}{@{}p{1.6cm}p{3.2cm}Lp{3.2cm}@{}}
\toprule
类型 & 内容 & 论文中常见用途 & 典型来源 \\
\midrule
公司文本 & 年报、MD\&A（管理层讨论与分析）、理财公告、ESG/社会责任报告 & ① 衡量企业风险、不确定性与战略取向；② 识别企业金融化、治理行为与信息披露质量；③ 构造 ESG、绿色、数字化等“软信息”指标 & 巨潮资讯网、交易所官网、公司官网 \\
媒体与舆情 & 新闻报道、投资者讨论（股吧）、搜索行为 & ① 衡量媒体关注度与舆论监督；② 刻画投资者情绪与信息不对称；③ 构造外部冲击或信息扩散指标 & 百度新闻、财经网站、股吧、百度指数 \\
平台与商业行为 & 电商销售、用户评价、招聘信息 & ① 反映市场需求与企业经营状况；② 衡量竞争强度、品牌影响力；③ 刻画用工需求、技术升级与数字化转型 & 淘宝/京东、大众点评、招聘网站 \\
政策与监管 & 政策文件、监管公告、环保处罚信息 & ① 构造政策强度或不确定性指标；② 作为\term{准自然实验}{quasi-natural experiment}的政策冲击变量；③ 衡量监管约束与合规风险 & 政府官网、部委网站、行政处罚公示平台 \\
\bottomrule
\end{tabularx}
\end{center}

\subsection{技术路线：从研究问题到数据来源}
网页结构 → 正则 → 百度新闻 → Selenium → IP 代理与反爬 → 清洗与文本过滤 → 巨潮资讯网 → PDF 解析（前置：Python 环境 + Chrome）。\key{对应}：公司文本 → Selenium 翻页下载 PDF + pdfplumber 解析 + os 批量处理；媒体舆情 → requests+headers、正则提取、改 URL 翻页、正文去噪；平台商业 → 同上，多为动态加载需 Selenium，反爬严需 IP 代理与限速；政策监管 → requests/Selenium 取列表与正文、PDF 解析。答题按\key{类型—内容—用途—来源}；设计题落到\key{网站、工具、字段、变量}（§\ref{sec:D06-design}）。

% ======================================================================
\section{课程考核要求（简述）}
\key{构成}：两次随堂小测各 30 分 + 课后作业 40 分；每次交\key{代码 + 运行视频}（从程序开始运行录到展示结果）；教学安排为 12 月 22 日至次年 1 月 19 日共 8 次（第 16–20 周）。\key{提交}：发指定作业邮箱，标题“姓名+学号”，\key{以邮件时间为准，超时视为未交、不可补交}；小测当天、课后作业下次课前一天 24:00 前交。\key{任务}：小测 1——四大行多页百度新闻（§\ref{sec:D06-quiz1}）；小测 2——阿里巴巴、工商银行 2 页新闻正文；课后作业——巨潮年报 PDF 批量下载，文件名含代码、年份、公司名（后两项见第 7 章）。

% ======================================================================
\section{网页源代码与开发者工具}
\begin{multicols}{2}\raggedcolumns\sloppy
\subsection{打开开发者工具：F12（笔记本 Fn+F12）}
装 Chrome（\texttt{google.cn/chrome}，火狐等亦可）。\key{F12} 打开开发者工具（DevTools）；笔记本按 \key{Fn+F12}，或右键 →“检查”。上方标签页依次为 Elements、Console、Sources、Network、Performance、Memory、Application、Security 等，爬虫只用 \key{Elements} 与 \key{Network}。

\subsection{Elements 面板（DOM 树）vs 查看网页源代码}
Elements 中的内容可理解为网站源码：\texttt{<html>}、\texttt{<head>}、\texttt{<body>}、\texttt{<div id="wrapper">} 层层嵌套，点小三角展开/折叠（看不到中文就逐层展开）。\tagext\ Elements 是执行 JavaScript 后的\term{文档对象模型}{Document Object Model, DOM}；“查看网页源代码”（Ctrl+U）是服务器最初返回的 HTML，即 \texttt{requests.get(url).text} 所得。\key{动态加载}网页（巨潮、新版百度）Elements 里有而 requests 可能拿不到 → 用 Selenium 取 \texttt{page\_source}。

\subsection{“选择元素”箭头（Ctrl+Shift+C）}
点左上角箭头（变蓝）→ 鼠标在页面移动，被指区域高亮，Elements 自动跳到对应源码；悬浮提示“标签名.class 名 + 尺寸”（如 \texttt{a.JNXyal}）。搜索关键词常被 \texttt{<em>} 或 \texttt{<font color="\#CC0000">} 包住标红。

\subsection{定位新闻标题与链接的三步法}
\key{① 点击}选择元素箭头 → \key{② 选择}页面上的新闻标题 → \key{③ 点击展开} \texttt{<a} 前的三角看完整内容。结构见 §\ref{sec:D06-box}：\texttt{href="..."} 里是\key{新闻链接}，\texttt{<a ...>} 与 \texttt{</a>} 之间是\key{新闻标题}（关键词被 \texttt{<em>} 包着）；提取规则都从这种观察出发。

\subsection{Network 面板：Headers、状态码、找 UA}
记录页面加载时的每个请求（\key{打开面板后需刷新页面}）；左侧请求列表（Name 列每个文件一行，过滤器 All/XHR/JS/CSS/Img/Doc，底部显示请求数与传输量，如“32 requests | 1.1 kB transferred”）。点任一请求 → Headers 三块：\key{General}（Request URL、Request Method: GET、Status Code: 200 OK、Referrer Policy）；\key{Response Headers}（如 Content-Type）；\key{Request Headers}（含 User-Agent）。\key{UA 查找路径}：F12 → Network → 刷新 → 点任一请求 → Headers → Request Headers → User-Agent。
\end{multicols}

% ======================================================================
\section{HTML 结构基础与“大框—中框—小框”定位思路}

\subsection{标签、属性、嵌套与注释：HTML 基本语法}
\term{超文本标记语言}{HyperText Markup Language, HTML}是网页的源代码语言。\key{标签}（tag）：尖括号标记，一般成对，\texttt{<p>} 开始、\texttt{</p>} 结束，中间是内容；少数单独出现（\texttt{<img ...>}、\texttt{<input ...>}）。\key{属性}（attribute）：开始标签里的“名=值”，如 \texttt{<div class="result" id="2">}；\key{id 通常全页唯一}，最适合定位；\key{class 可被多个元素共用}，一个元素可有多个 class（空格分隔：\texttt{class="c-summary c-row "}）；\texttt{data-} 开头（\texttt{data-click}、\texttt{data-id}、巨潮的 \texttt{data-seccode}）为网站自定义数据。\key{嵌套}（nesting）：标签套标签形成 DOM 树，目标信息几乎总在\key{最内层}。注释 \texttt{<!-{}- ... -{}->} 不显示。

\subsection{“大框—中框—小框”：一条百度新闻结果的嵌套结构}\label{sec:D06-box}
\begin{lstlisting}[style=txtD06]
<div class="result" id="2">                          <- 大框：一条新闻（各条 id="2"、"3"…依次排列）
    <h3 class="c-title">                             <- 中框：标题区
      <a href="http://finance.sina.com.cn/roll/2018-12-18/doc-..."
         data-click="{'f0':'77A717EA', ...}" target="_blank">        <- 小框：链接 + 标题
         <em>阿里巴巴</em>马云:创新是解决问题的唯一途径</a></h3>
    <div class="c-summary c-row ">...</div>         <- 摘要、来源、日期
</div>
\end{lstlisting}
要诀：\key{内容一般都在最后的小框里}；\key{点击箭头，展开大框，看到小框}。分框是为了找到每条都相同、又能唯一圈住目标的\key{左、右边界}：直接搜 \texttt{<a href="} 会抓到导航、广告等所有链接；先锁定中框 \texttt{<h3 class="c-title">} 再找 \texttt{<a href="(...)"} 只命中新闻——即“文本A\texttt{(.*?)}文本B”（§\ref{sec:D06-AB}）。

\subsection{常用标签速查：div、h1–h6、p、a（target="\_blank"）及其他}
\begin{center}\footnotesize\renewcommand{\arraystretch}{1.0}
\begin{tabularx}{\linewidth}{@{}p{3.3cm}L@{}}
\toprule
\texttt{<div>} & “空盒子”（division），本身无样式含义；外层 div 多用于\key{布局}（圈出一条新闻、导航栏、正文区），靠 class/id 区分；正文抽取靠正文 div 的 id/class 或“文本最长的 div” \\
\texttt{<h1>}–\texttt{<h6>} & 标题，\texttt{<h1>标题内容</h1>}；\key{h1 字号最大、h6 最小}；百度新闻标题在 \texttt{<h3>} \\
\texttt{<p>} & 段落，\texttt{<p>段落内容</p>}；新闻正文由一串 \texttt{<p>} 组成；百度来源与日期在 \texttt{<p class="c-author">} \\
\texttt{<a>} & 链接，\texttt{<a href="链接地址">文本内容</a>}：\texttt{href} 为跳转网址，标签间是显示文字；加 \texttt{target=\_blank}（规范写 \texttt{target="\_blank"}）则在\key{新标签页}打开、不覆盖原页 \\
\midrule
\tagext\ 其他标签 & \texttt{<!DOCTYPE html>} 文档声明，\texttt{<html>} 根标签；\texttt{<head>} 头部（不显示）、\texttt{<body>} 主体；\texttt{<em>}、\texttt{<font>} 强调/标红搜索关键词（提取后 \texttt{re.sub('<.*?>','',s)} 去掉）；\texttt{<span>} 行内小容器（巨潮公告标题 \texttt{<span class="r-title">}）；\texttt{<table>}／\texttt{<tr>}／\texttt{<td>} 表格／行／单元格（巨潮结果行 \texttt{<tr class="el-table\_\_row">}）；\texttt{<input>}、\texttt{<button>} 输入框、按钮（百度搜索框 \texttt{<input id="kw">}）；\texttt{<img>} 图片；\texttt{<script>} 脚本 \\
\tagext\ HTML 实体 & \texttt{\&nbsp;} = 不换行空格（百度用 \texttt{\&nbsp;\&nbsp;} 隔开来源与日期）；字符 \texttt{\&} 写作 \texttt{\&amp;} \\
\bottomrule
\end{tabularx}
\end{center}

\subsection{最简单的网页 test.html}
\begin{lstlisting}[style=txtD06]
<!DOCTYPE html>
<html><body>
    <h1>这是标题 1</h1>  <p>这是标题1下的段落。</p>  <h2>这是标题 2</h2>
    <a href="https://www.baidu.com" target=_blank>这是带链接的内容</a>
</body></html>
\end{lstlisting}
存为 test.html 用浏览器打开：大号标题 1、一段正文、稍小标题 2、一个可点击链接（\texttt{target=\_blank}：新标签页打开百度）。

% ======================================================================
\section{HTTP 请求、URL 与 User-Agent}

\subsection{HTTP 请求—响应模型：浏览器与爬虫做同一件事}
客户端（浏览器/程序）发 \key{HTTP 请求}（GET + URL + headers）→ 服务器回 \key{HTTP 响应}（状态码 + 响应头 + HTML）。浏览器把源码“画”成页面；爬虫把源码当\key{字符串}交给程序处理。

\subsection{URL 结构：协议、域名、路径、查询参数（word、pn）}
\term{统一资源定位符}{Uniform Resource Locator, URL}，例：\url{https://www.baidu.com/s?rtt=1&bsst=1&cl=2&tn=news&ie=utf-8&word=阿里巴巴&pn=10}
\begin{center}\footnotesize\renewcommand{\arraystretch}{1.0}
\begin{tabularx}{\linewidth}{@{}p{2.5cm}Lp{2.6cm}L@{}}
\toprule
\texttt{https://} & 协议（http／加密的 https）& \texttt{tn=news} & 搜索类型：新闻 \\
\texttt{www.baidu.com} & 域名（哪台服务器）& \texttt{ie=utf-8} & 输入编码 \\
\texttt{/s} & 路径（服务器上的页面/程序）& \texttt{word=阿里巴巴} & 关键词——\key{批量爬多家公司只换它} \\
\texttt{?} 之后 & 查询参数（query parameters），“名=值”以 \texttt{\&} 连接 & \texttt{pn=10} & 偏移量 0/10/20 = 第 1/2/3 页，\key{翻页只改它} \\
& & \texttt{rtt=1\&bsst=1\&cl=2} & 百度内部排序/显示参数，照抄 \\
\bottomrule
\end{tabularx}
\end{center}
汉字传输时转成百分号编码（\texttt{\%E9\%98\%BF...}）。\key{“URL 找不同”}：\texttt{url = '...\&word=' + company + '\&pn=' + str((page-1)*10)}。

\begin{multicols}{2}\raggedcolumns\sloppy
\subsection{GET/POST 与状态码 200/30x/403/404/500}
\key{GET}：取网页，参数写在 URL 里（本章所有请求）；\tagext\ \key{POST}：提交表单。状态码：\key{200 OK} 成功；\tagext\ 301/302 重定向（requests 默认自动跟随）；\key{403 Forbidden} 禁止访问（没加 User-Agent、IP 被封）；\tagext\ 404 Not Found 页面不存在（网址拼错、链接失效）；\tagext\ 500 服务器内部错误（对方出错，稍后重试）。

\subsection{headers 与 User-Agent：“身份证”}
\term{用户代理}{User-Agent, UA} 告诉服务器“你是谁——浏览器？程序？爬虫？”；requests 默认 UA 为 \texttt{python-requests/版本号}，一眼可认出。\key{不加 UA 的三种后果}：① 返回空页面；② 直接 403；③ 返回“简化版页面”（内容不全）。做法：从 Network 复制浏览器 UA，放进\key{字典} \texttt{headers = \{'User-Agent': 'Mozilla/5.0 ...'\}}，以 \texttt{requests.get(url, headers=headers)} 随请求发出，把程序“伪装”成浏览器（完整字符串见 §\ref{sec:D06-full}，表示 Windows 10 64 位上的 Chrome 69）。\textbf{易错}：键名写成 \texttt{'User\_Agent'}/\texttt{'UserAgent'}（不被识别，等于没加）；把 UA 字符串直接传给 \texttt{headers=}（必须是字典）。\tagext\ 有的网站还查 \texttt{Referer}（来源页）、\texttt{Cookie}（登录状态），同样复制进字典。
\end{multicols}

% ======================================================================
\section{requests 库：发请求、拿源码}

\subsection{requests 是什么：HTTP 客户端库}
\key{requests} 是 HTTP 客户端库（HTTP client library），本质一句话：\key{向服务器发请求，拿响应结果}。非标准库，首次使用需安装：Win+R → cmd → \texttt{pip install requests}（Anaconda 已自带）。

\subsection{最小程序：requests.get() 发请求，.text 取字符串 HTML}\label{sec:D06-min}
\begin{lstlisting}[style=pyD06]
import requests
headers = {'User-Agent': 'Mozilla/5.0 (Windows NT 10.0; Win64; x64) ...'}           # 伪装浏览器，完整 UA 见案例完整程序
url = 'https://www.baidu.com/s?rtt=1&bsst=1&cl=2&tn=news&ie=utf-8&word=阿里巴巴'  # 浏览器中复制，汉字可直接写
res = requests.get(url, headers=headers).text; print(res)  # get() 返回响应对象；.text 取源码字符串
\end{lstlisting}
\key{背后过程}：GET 请求（带 headers）→ 服务器返回响应（状态码、响应头、HTML）→ 响应对象 → \texttt{.text} 解码为字符串。\textbf{易错}：\texttt{requests.get(...)} 本身是\key{响应对象}（打印为 \texttt{<Response [200]>}），不是字符串；漏写 \texttt{.text} 时 \texttt{'关键词' in res} 不报错但恒为 False，\texttt{re.findall(p, res)} 报 \texttt{TypeError: expected string or bytes-like object}。

\begin{multicols}{2}\raggedcolumns\sloppy
\subsection{检验是否拿到网页：Ctrl+F 与 in}
在输出（或存成文件的 res）中 Ctrl+F 搜“阿里巴巴”：找得到标题 → 拿到内容；找不到或源码很短 → \key{被拦截}（没加 UA、触发验证码）或\key{动态加载}（需 Selenium）；程序内用 \texttt{print('阿里巴巴' in res)}。

\subsection{响应对象的常用属性 \texorpdfstring{\tagext}{（拓展）}}
\texttt{r = requests.get(url, headers=headers)}：\texttt{r.text} 按 \texttt{r.encoding} 解码的字符串源码；\texttt{r.content} 原始字节（下载 PDF、图片写文件用）；\texttt{r.status\_code} 状态码；\texttt{r.encoding} 当前解码编码，\texttt{r.apparent\_encoding} 按内容推测的编码（中文乱码时先 \texttt{r.encoding = r.apparent\_encoding} 再取 \texttt{r.text}）；\texttt{r.url} 最终网址（汉字已百分号编码）。

\subsection{get 参数：headers/timeout/proxies/params}
\texttt{requests.get(url, headers=headers, timeout=10, proxies=proxies).text}：\texttt{headers} 请求头字典（伪装浏览器）；\texttt{timeout=10} 10 秒无响应即抛异常、防卡死；\texttt{proxies} IP 代理字典（第 7 章）。\tagext\ \texttt{params=\{'tn':'news', 'word':'阿里巴巴', 'pn':10\}} 由 requests 自动拼接并编码查询参数。
\end{multicols}

% ======================================================================
\section{正则表达式基础：findall 与元字符}

\subsection{为什么用正则：从一长串源码中“抠”出内容}
源码是几十万字符的字符串，目标（标题、链接、来源、日期）总在固定标签之间、格式有规律。\term{正则表达式}{regular expression, regex}是描述字符串格式规律的小语言：用\key{模式}（pattern）匹配文本，找出符合规则的片段。标准库 \texttt{re}，\texttt{import re} 即可（无需安装）。

\subsection{re.findall(pattern, string, flags=0) 的返回规则}
找出\key{所有不重叠}的匹配，\key{结果总是列表}：多个匹配 → 多个元素（\lstinline|re.findall(r'\d\d\d', 'Hello 123 world 456')| → \texttt{['123', '456']}）；\key{无匹配 → \texttt{[]}}（不报错、不是 None）；规则\key{无括号}返回整个匹配，\key{一个括号}只返回括号内内容，\key{多个括号}返回元组列表（§\ref{sec:D06-tuple}）；\texttt{flags} 为修饰符，最常用 \texttt{re.S}（§\ref{sec:D06-reS}）。

\subsection{\texorpdfstring{原始字符串 \texttt{r'...'} 与转义问题}{原始字符串 r'...' 与转义问题}}
反斜杠在 Python 字符串和正则里各有一层含义：普通字符串中 \lstinline|'\n'| 换行、\lstinline|'\t'| 制表符、\lstinline|'\\'| 一个反斜杠；\key{原始字符串}（raw string，引号前加 \texttt{r}）中反斜杠不转义，\lstinline|r'\d'| 即“\textbackslash”“d”两个字符，原样交给正则引擎。\lstinline|'\d'| 是非法转义：旧版 Python 照常保留，\key{3.12 起}报 \texttt{SyntaxWarning: invalid escape sequence}（仍能运行）。\key{正则一律加 r}。Windows 路径 \lstinline|'C:\Users\...'| 中 \lstinline|\U| 是 Unicode 转义开头，不加 r 直接 \texttt{SyntaxError}，应写 \lstinline|r'C:\Users\...'|。

\subsection{常用元字符速查表}
注意“一个”：\texttt{\textbackslash d}、\texttt{\textbackslash w}、\texttt{.} 等每次只匹配\key{一个字符}；例子见下列各小节。
\begin{center}\footnotesize\renewcommand{\arraystretch}{1.0}
\begin{tabularx}{\linewidth}{@{}p{1.3cm}Lp{1.6cm}L@{}}
\toprule
模式 & 描述 & 模式 & 描述 \\
\midrule
\texttt{\textbackslash d} & 一个数字（0–9；Unicode 下还含全角数字等）& \texttt{*} & 前一个单位重复 \key{0 次或多次} \\
\texttt{\textbackslash w} & 一个字母、数字或下划线（Python 3 中\key{也匹配汉字}）& \texttt{+} & 前一个单位重复 \key{1 次或多次} \\
\texttt{\textbackslash s} & 一个空白字符（空格、\texttt{\textbackslash t}、\texttt{\textbackslash n} 等）& \texttt{?} & 常与 \texttt{.}、\texttt{*} 组成\key{非贪婪}匹配 \texttt{.*?}（尽量少匹配）\\
\texttt{\textbackslash S} & 一个非空白字符 & \texttt{( )} & 匹配括号内表达式，也表示一个组 \\
\texttt{\textbackslash n}／\texttt{\textbackslash t} & 一个换行符／一个制表符（Tab）& \texttt{.} & 一个任意字符，\key{换行符除外} \\
\midrule
\multicolumn{4}{@{}l}{\tagext}\\
\texttt{?}（单独）& 前一个单位出现 0 或 1 次 & \texttt{[\^{} ]} & 取反：不在方括号中的任一字符 \\
\texttt{\{m,n\}} & 重复 m 到 n 次；\texttt{\{m\}} 恰好 m 次；\texttt{\{m,\}} 至少 m 次 & \texttt{|} & 或，如 \texttt{r'(五|5|主要)'} \\
\texttt{[ ]} & 字符集：任选其一；其中元字符失去特殊含义 & \texttt{\^{}}、\texttt{\$} & 字符串开头、结尾（不占字符）\\
& & \texttt{\textbackslash.}\,\texttt{\textbackslash*}\,\texttt{\textbackslash(} & 转义：匹配字面的点、星号、括号（\texttt{r'\textbackslash.'} 对 \texttt{'3.14'} → \texttt{['.']}）\\
\bottomrule
\end{tabularx}
\end{center}

\begin{multicols}{2}\raggedcolumns\sloppy
\subsubsection{\texorpdfstring{\texttt{\textbackslash d}：一个数字；\texttt{\textbackslash d\textbackslash d\textbackslash d} 与 \texttt{\textbackslash d+} 的区别}{\textbackslash d、\textbackslash d\textbackslash d\textbackslash d 与 \textbackslash d+ 的区别}}
对 \texttt{'订单号A12345'}：\lstinline|r'\d'| → \texttt{['1','2','3','4','5']}（每次吃一个，逐个找出）；\lstinline|r'\d\d\d'| → \texttt{['123']}（45 不够三位）；\lstinline|r'\d+'| → \texttt{['12345']}。\textbf{易错}：6 位股票代码用 \texttt{\textbackslash d\{6\}}，\texttt{\textbackslash d+} 会把价格 \texttt{1650.5} 拆成 \texttt{'1650'}、\texttt{'5'} 一并抓出。

\subsubsection{\texorpdfstring{\texttt{\textbackslash w}：单词字符，在 Python 3 中也匹配中文}{\textbackslash w：单词字符，也匹配中文}}
\lstinline|r'\w+'| 对 \texttt{'abc\_2025'} → \texttt{['abc\_2025']}。Python 3 的 re 默认 Unicode 规则：\lstinline|re.findall(r'\w+', '阿里巴巴2025')| → \texttt{['阿里巴巴2025']}。只要英文字母、数字、下划线：\lstinline|r'[A-Za-z0-9_]+'| → \texttt{['2025']}；或 \tagext\ \texttt{re.ASCII} 模式：\lstinline|re.findall(r'\w+', '阿里巴巴2025', re.ASCII)| → \texttt{['2025']}。

\subsubsection{\texorpdfstring{\texttt{\textbackslash s}／\texttt{\textbackslash S}、\texttt{\textbackslash n}／\texttt{\textbackslash t}：空白、非空白、换行、制表符}{\textbackslash s、\textbackslash S、\textbackslash n、\textbackslash t：空白、非空白、换行、制表符}}
\lstinline|re.findall(r'\s', 'a b\tc\nd')| → \texttt{[' ', '\textbackslash t', '\textbackslash n']}；\lstinline|re.findall(r'\S+', ' 中南 大学 2025 ')| → \texttt{['中南','大学','2025']}；\tagext\ \lstinline|re.sub(r'\s+', ' ', s)| 把连续空白压成一个空格。\texttt{\textbackslash n} 匹配换行符（\lstinline|'第一行\n第二行'| → \texttt{['\textbackslash n']}），\texttt{\textbackslash t} 匹配 Tab（\lstinline|'姓名\t学号'|）；源码与 PDF 文本充满换行 → 需 \texttt{re.S}。

\subsubsection{点号 .：占一个位置，什么字符都行，但只能一个}
\lstinline|r'c.t'| 对 \texttt{'cat cut cot'} → \texttt{['cat','cut','cot']}。易错：① 只是\key{一个}字符（\texttt{ct} 不匹配；空格可匹配）；② \key{不跨行}：\lstinline|re.findall(r'c.t', 'c\nt')| → \texttt{[]}；③ 一段任意字符用 \texttt{.*}。

\subsubsection{* 与 +：重复前一个单位（作用范围）}
\lstinline|r'ab*c'| 对 \texttt{'ac abc abbbc adc'} → \texttt{['ac','abc','abbbc']}；\lstinline|r'ab+c'| 对 \texttt{'ac abc abbbc'} → \texttt{['abc','abbbc']}；\tagext\ \lstinline|r'ab?c'| → \texttt{['ac','abc']}。\texttt{*} 只作用于\key{紧挨在前的一个单位}（b 而非 ab），整体重复用 \texttt{(ab)*}。故 \texttt{.*} = 一段任意长度的任意字符（不跨行）。

\subsubsection{\texorpdfstring{量词 \texttt{\{m,n\}}：限定重复次数 \tagext}{量词 \{m,n\}：限定重复次数（拓展）}}
对 \texttt{'2025年12月18日'}：\lstinline|r'\d{4}'| → \texttt{['2025']}；\lstinline|r'\d{1,2}'| → \texttt{['20','25','12','18']}（贪婪，尽量取 2 位）；日期 \lstinline|r'\d{4}-\d{2}-\d{2}'| → \texttt{['2025-12-18']}。“前五大客户”正则 \texttt{(五|5|主要).\{0,10\}?客户}：\texttt{.\{0,10\}?} = 0–10 个任意字符、非贪婪（第 7 章）。

\subsubsection{\texorpdfstring{字符集 \texttt{[ ]}、取反 \texttt{[\^{} ]} 与“或” \texttt{|} \tagext}{字符集 [ ]、取反与“或”（拓展）}}
\texttt{[abc]} 任选一个；\texttt{[0-9]} 只含半角数字（\texttt{\textbackslash d} 还匹配全角“１２”）；\texttt{[A-Za-z0-9\_]} 为英文单词字符；常用汉字 \texttt{[一-龥]} ≡ \lstinline|[\u4e00-\u9fa5]|（Unicode 4E00–9FA5，统计汉字个数用）。方括号内开头的 \texttt{\^{}} 表示取反（\lstinline|r'[^0-9]+'| 对 \texttt{'A1B22'} → \texttt{['A','B']}）；方括号中 \texttt{*}、\texttt{.} 失去特殊含义（\texttt{'[*]'} 的原理，§\ref{sec:D06-star}）。\texttt{A|B} 表示 A 或 B，常配括号限定范围：\lstinline!re.findall(r'(工商|建设)银行', '工商银行和建设银行')! → \texttt{['工商','建设']}（有括号只返回括号内容）。

\subsubsection{\texorpdfstring{锚点 \texttt{\^{}} 与 \texttt{\$}：开头与结尾 \tagext}{锚点 \textasciicircum{} 与 \$：开头与结尾（拓展）}}
\lstinline|re.findall(r'^\d+', '123abc456')| → \texttt{['123']}；\lstinline|re.findall(r'\d+$', '123abc456')| → \texttt{['456']}。\texttt{\$} 可为 \texttt{(.*?)} 充当右边界（§\ref{sec:D06-empty}）。
\end{multicols}

% ======================================================================
\section{非贪婪匹配 (.*?) 与分组}

\subsection{贪婪 .* 与非贪婪 .*?：\texorpdfstring{\texttt{<title>}}{<title>} 例与匹配机制}\label{sec:D06-greedy}
\begin{lstlisting}[style=pyD06]
t = '<title>AAA</title><title>BBB</title>'
re.findall('<title>(.*)</title>', t)    # ['AAA</title><title>BBB']  贪婪：吃太多，只 1 个元素
re.findall('<title>(.*?)</title>', t)   # ['AAA', 'BBB']             非贪婪：刚好一段
\end{lstlisting}
\term{贪婪匹配}{greedy matching} \texttt{.*}：在能匹配成功的前提下\key{尽量多吃}——先吃到末尾，再逐字符往回“吐”，直到后面接得上 \texttt{</title>}，第一次接上的是\key{最后一个} \texttt{</title>}，中间的 \texttt{</title><title>} 全被吞掉。\term{非贪婪匹配}{non-greedy / lazy matching} \texttt{.*?}：先吃 0 个，不行再多吃一个，遇到\key{第一个} \texttt{</title>} 就停。源码中同样的标签成百上千次重复，爬虫几乎总用 \texttt{.*?}。\textit{Greedy quantifiers match as much as possible; adding a question mark makes them non-greedy.}

\begin{multicols}{2}\raggedcolumns\sloppy
\subsection{分组 ()：整体重复、捕获；非捕获组 (?:)}\label{sec:D06-group}
括号两个作用：① \key{整体}：\texttt{(ab)+} 匹配 ab、abab、ababab；② \key{捕获}（capturing group，最常用）：findall 遇到带括号的规则\key{只返回括号里的内容}，两边“定位文字”不返回。\tagext\ 重复的组只记住\key{最后一次}：\lstinline|re.findall('(ab)+', 'ab abab ababab')| → \texttt{['ab','ab','ab']}；要整段用非捕获组 \lstinline|(?:ab)+| → \texttt{['ab','abab','ababab']}。
\end{multicols}

\subsection{(.*?) 无右边界为什么得到空串（“日期：(.*?)”）}\label{sec:D06-empty}
\begin{lstlisting}[style=pyD06]
s2 = '来源：新华社 日期：2025-12-18'
re.findall(r"日期：(.*?)", s2)               # ['']  无右边界：吃 0 个就满足
# 正确（均得 ['2025-12-18']）：r"日期：(.*)" 贪婪到行尾；r"日期：(\d{4}-\d{2}-\d{2})"；r"日期：(.*?)$"
re.findall(r"来源：(.*?) 日期", s2)          # ['新华社']  夹在两段文字之间
\end{lstlisting}
\texttt{(.*?)} 能少吃就少吃，无右边界时吃 0 个就满足 → \texttt{['']}。\key{原则：\texttt{(.*?)} 必须夹在文本 A 与文本 B 之间}（找错高频）。

\begin{multicols}{2}\raggedcolumns\sloppy
\subsection{核心套路：文本A(.*?)文本B}\label{sec:D06-AB}
\begin{thm}[“文本A(.*?)文本B”提取法]
目标 X 左侧紧挨、每处相同的文字作\key{文本 A}，右侧紧挨的相同文字作\key{文本 B}，则 \texttt{re.findall('文本A(.*?)文本B', 源码, re.S)} 返回所有 X 的列表。文本 A 要\key{足够独特}（如带 class 的开始标签），文本 B 通常取\key{结束标签}（\texttt{</a>}、\texttt{"}）。
\end{thm}
\lstinline|re.findall('文本A(.*?)文本B', '文本A百度新闻文本B')| → \texttt{['百度新闻']}：不确定内容的长度与格式，但知道它在哪两块内容中间。\textbf{易错}：\texttt{'<span.*?>(.*?)</span>'} 会把价格、名称两个 span 都抓到无法区分；应分别写 \texttt{'<span class="price">(.*?)</span>'}、\texttt{'<span class="name">(.*?)</span>'}。

\subsection{多个匹配与规则变量 p\_xxx}
res 中有三段“文本A…文本B”，分别夹着百度新闻、新浪财经、搜狐新闻：\texttt{p\_source = '文本A(.*?)文本B'}，\texttt{re.findall(p\_source, res)} → \texttt{['百度新闻', '新浪财经', '搜狐新闻']}（依次全部找出）；误用贪婪 \texttt{(.*)} 只得 1 个元素（从第一个文本A 吃到最后一个文本B）。\key{p = pattern}：规则单独存成变量（\texttt{p\_title}、\texttt{p\_href}），便于修改、复用。

\subsection{.*? 填充 vs (.*?) 提取}
\texttt{'<h3 class="c-title">.*?>(.*?)</a>'}：\texttt{<h3 class="c-title">.*?>} 整体是文本 A，\texttt{</a>} 是文本 B；\key{不带括号的 \texttt{.*?} 填充}不关心、每条不同的内容（\texttt{<a} 后的 href、data-click、target），跳过不返回；\key{带括号的 \texttt{(.*?)} 提取}目标（标题）。简化例：\lstinline|re.findall('<h3>文本C.*?文本D(.*?)</h3>', '<h3>文本C<变化的网址>文本D新闻标题</h3>')| → \texttt{['新闻标题']}（每条都不同的“变化的网址”无法写进规则，用 \texttt{.*?} 跳过）。
\end{multicols}

\subsection{多个分组：findall 返回元组列表 \texorpdfstring{\tagext}{（拓展）}}\label{sec:D06-tuple}
元组按括号顺序装各组内容；相比每字段单独一条规则（等长列表按下标对应），不会因某字段少抓一条而错位：
\begin{lstlisting}[style=pyD06]
s = '<a href="1.html" target="_blank">标题一</a><a href="2.html"><em>阿里</em>标题二</a>'
re.findall('<a href="(.*?)".*?>(.*?)</a>', s)  # [('1.html', '标题一'), ('2.html', '<em>阿里</em>标题二')]
\end{lstlisting}

% ======================================================================
\section{修饰符 re.S、替换与转义}

\subsection{修饰符 re.S（DOTALL）：让点号匹配换行}\label{sec:D06-reS}
点号 \texttt{.} 默认不匹配 \texttt{\textbackslash n}；\texttt{re.S}（DOTALL）使 \texttt{.} 也匹配换行，\texttt{.*?} 可\key{跨行}（“自动考虑换行的影响”）。
\begin{lstlisting}[style=pyD06]
res = '''文本A
       百度新闻文本B'''                            # 三引号字符串：文本A 后有换行和空格
re.findall('文本A(.*?)文本B', res)                 # []                   不加 re.S：跨不过换行
source = re.findall('文本A(.*?)文本B', res, re.S)  # ['\n       百度新闻']  带换行和空格
source[0].strip()                                  # '百度新闻'
\end{lstlisting}
\key{爬网页时 findall 几乎总加 re.S}。\textbf{易错}：目标与边界同在一行时不加 re.S 也能匹配；贪婪 \texttt{(.*)} 不加 re.S 只吃到本行，加了会跨行吃到最后一个右边界。

\begin{multicols}{2}\raggedcolumns\sloppy
\subsection{strip()：去掉提取结果两端的空白}
\texttt{字符串.strip()} 删除\key{两端}的空格、\texttt{\textbackslash n}、\texttt{\textbackslash t}，中间不动；“\key{re.S + strip()}”总成对出现（\texttt{lstrip}、\texttt{rstrip}、带参数用法见第 7 章）。

\subsection{str.replace(旧, 新)：固定字符串替换}
百度用 \texttt{<em>...</em>} 包住关键词，提取的标题带标签：\texttt{title[0] = title[0].replace('<em>', '')}，再 \texttt{.replace('</em>', '')}。replace 是字符串方法，一次只换\key{一个固定串}；字符串不可变、返回新串，\key{必须赋值回去}。

\subsection{re.sub(pattern, repl, string)：去掉所有 <...> 标签}
把 string 中所有匹配 pattern 的片段换成 repl，返回新串；不仅去掉 \texttt{<em></em>}，还去掉 \texttt{<font>}、\texttt{<img>} 等所有 \texttt{<>} 内容。\texttt{title[0] = re.sub('<.*?>', '', title[0])}：左 \texttt{title[0]} 为替换后标题；\texttt{'<.*?>'} = 所有“<任何内容>”（\texttt{.*?} 起填充作用）；\texttt{''} = 替换成空（删除）；右 \texttt{title[0]} 为原标题。\textbf{易错}：必须非贪婪；\texttt{'<.*>'} 从第一个 \texttt{<} 删到最后一个 \texttt{>}：\texttt{'<em>阿里巴巴</em>马云:…'} → \texttt{'马云:…'}。编辑器显示的 \texttt{pattern:}、\texttt{repl:} 只是参数名提示。
\key{对比}：\texttt{strip()} 去两端空白；\texttt{replace(old, new)} 换一个固定子串（所有出现处）；\texttt{re.sub(p, repl, s)} 按正则换所有匹配（去标签、去星号、压空白）；\texttt{split(sep)} 按分隔符切成列表。
\end{multicols}

\subsection{中括号 [ ] 与反斜杠转义：处理 *ST 公司名中的星号}\label{sec:D06-star}
中括号使其中内容\key{不再有特殊含义}。\texttt{*ST} 公司（被实施退市风险警示）名称含 \texttt{*}，而 \texttt{*} 在正则中表示“重复”：
\begin{lstlisting}[style=pyD06]
re.sub('[*]', '', '*阿里巴巴')   # '阿里巴巴'   [*] 表示字面星号
re.sub(r'\*', '', '*ST阿里')     # 'ST阿里'     另一写法：反斜杠转义
re.sub('*', '', '*阿里巴巴')     # 报错 re.error: nothing to repeat at position 0
\end{lstlisting}
\texttt{'*'} 前面无可重复的东西而报错（3.13 起类名显示为 \texttt{PatternError}）。\texttt{. ? + ( ) [ ] \{ \} | \^{} \$ \textbackslash} 同理；\tagext\ \texttt{re.escape('*ST康美(600518)')} 自动加反斜杠。公司名先去星号，再拼网址、正则规则（否则报错）或 Windows 文件名（不许含 \texttt{*}）。

\subsection{re 模块的其他函数：search、match、compile、split \texorpdfstring{\tagext}{（拓展）}}
\begin{center}\footnotesize\renewcommand{\arraystretch}{1.0}
\begin{tabularx}{\linewidth}{@{}p{3cm}L@{}}
\toprule
\texttt{re.findall(p, s)} & 所有匹配 → 列表（无匹配为 \texttt{[]}）\\
\texttt{re.search(p, s)} & 全串找\key{第一个}匹配 → Match 对象（找不到 → \texttt{None}，可写 \texttt{if re.search(p, s):}）；\texttt{re.search(r'\textbackslash d+', '订单号A12345')} 显示 \texttt{<re.Match object; span=(4, 9), match='12345'>}，\texttt{.group()} → \texttt{'12345'} \\
\texttt{re.match(p, s)} & \leavevmode\key{只从开头}匹配，开头不符即 \texttt{None}：\texttt{re.match(r'\textbackslash d+', '订单号A12345')} → \texttt{None} \\
\texttt{re.sub(p, repl, s)} & 替换所有匹配 → 新字符串 \\
\texttt{re.split(p, s)} & 按规则切分 → 列表：\texttt{re.split(r'[，,]', 'a，b,c')} → \texttt{['a','b','c']} \\
\texttt{re.compile(p)} & 预编译为正则对象 q，再 \texttt{q.findall(s)}／\texttt{q.search(s)}／\texttt{q.sub(...)}；循环中反复用更快、更整洁 \\
\bottomrule
\end{tabularx}
\end{center}

% ======================================================================
\section{案例实战：百度新闻爬取}

\subsection{任务与流程总览：完整程序（多家公司 × 多页）}\label{sec:D06-full}
目标：每条新闻的\key{标题、网址、来源、日期}：拼网址 → 取源码 → 正则提取 → 清洗 → 输出，外套“公司 × 页码”双层循环（分步见 §\ref{sec:D06-step12}–§\ref{sec:D06-pn}）。
\begin{lstlisting}[style=pyD06]
import requests, re, time
headers = {'User-Agent': 'Mozilla/5.0 (Windows NT 10.0; Win64; x64) AppleWebKit/537.36 (KHTML, like Gecko) Chrome/69.0.3497.100 Safari/537.36'}
def baidu(company, page):
    num = (page - 1) * 10                                                    # 1 拼网址：pn=(页码-1)*10
    url = 'https://www.baidu.com/s?rtt=1&bsst=1&cl=2&tn=news&ie=utf-8&word=' + company + '&pn=' + str(num)
    res = requests.get(url, headers=headers, timeout=10).text                # 2 取源码字符串
    href  = re.findall('<h3 class="c-title">.*?<a href="(.*?)"', res, re.S)  # 3 提取：链接
    title = re.findall('<h3 class="c-title">.*?>(.*?)</a>', res, re.S)       #   标题（带 <em> 和空白）
    info  = re.findall('<p class="c-author">(.*?)</p>', res, re.S)           #   “来源&nbsp;&nbsp;日期”
    source, date = [], []                                                    # 4 清洗：来源与日期
    for i in range(len(info)):
        info[i] = re.sub('<.*?>', '', info[i])                               #   去 <img ...> 等标签
        source.append(info[i].split('&nbsp;&nbsp;')[0].strip())
        date.append(info[i].split('&nbsp;&nbsp;')[1].strip())
    with open('数据挖掘报告.txt', 'a', encoding='utf-8') as f:               # 5 输出：追加写入 txt
        f.write(company + '第' + str(page) + '页\n')
        for i in range(len(title)):
            title[i] = re.sub('<.*?>', '', title[i]).strip()                 #   清洗标题
            f.write(str(i + 1) + '.' + title[i] + '(' + date[i] + '-' + source[i] + ')\n')
            f.write(href[i] + '\n')
companys = ['工商银行', '农业银行', '中国银行', '建设银行']
for company in companys:              # 外层：公司
    for page in range(1, 6):          # 内层：第 1-5 页
        try:
            baidu(company, page); print(company + '第' + str(page) + '页爬取成功')
        except:
            print(company + '第' + str(page) + '页爬取失败')
        time.sleep(1)                 # 每次请求间隔 1 秒，降低被封概率
\end{lstlisting}

\subsection{第一步与第二步：取源码，写出标题、链接、来源与日期的正则}\label{sec:D06-step12}
F12 观察（§\ref{sec:D06-box}）：标题、链接在中框 \texttt{<h3 class="c-title">} 内的 \texttt{<a>}；来源、日期在 \texttt{<p class="c-author">来源\&nbsp;\&nbsp;日期</p>}。\texttt{p\_href} 用中框开头锁定，\texttt{(.*?)} 夹在 \texttt{<a href="} 与下一个 \texttt{"} 之间；\texttt{p\_title} 中 \texttt{.*?>} 跳过 \texttt{<a ...>} 属性串到第一个 \texttt{>}，\texttt{</a>} 为文本 B；\texttt{p\_info} 先整段取出来源与日期。三条都加 \texttt{re.S}（\texttt{<h3>} 与 \texttt{<a>} 间有换行）；三个列表\key{按下标一一对应}（下标 0 为第 1 条）。

\begin{multicols}{2}\raggedcolumns\sloppy
\subsection{第三步：清理标题（strip + re.sub）}
\texttt{for i in range(len(title)):}（下标 $0,\dots,n-1$）中 \texttt{title[i] = title[i].strip()} 去两端换行空格，再 \texttt{title[i] = re.sub('<.*?>', '', title[i])} 去 \texttt{<em>} 等标签。\textbf{易错}：须用下标原地修改，\texttt{for t in title: t = t.strip()} 只改临时变量；先 strip 或先 sub 均可，标签前后夹空白时 sub 后再 strip 最保险。

\subsection{第四步：清理来源与日期（去标签→split→strip）}
\begin{sloppypar}info 元素形如 \texttt{'新浪财经\&nbsp;\&nbsp;2025年12月18日 10:30'}（有时前带媒体 logo 的 \texttt{<img ...>}）。建 \texttt{source, date = [], []}，循环中 \texttt{re.sub('<.*?>', '', info[i])} 去标签 → \texttt{split('\&nbsp;\&nbsp;')} 切开，\texttt{[0]} 来源、\texttt{[1]} 日期，各 \texttt{.strip()} 后 \texttt{append}。日期有“2025年12月18日 10:30”“3小时前”“昨天 20:26”等写法，统一为 YYYY-MM-DD 见第 7 章。\end{sloppypar}

\subsection{第五步：打印输出与 str() 拼接}
\begin{sloppypar}\texttt{print(str(i + 1) + '.' + title[i] + '(' + date[i] + '-' + source[i] + ')')}，下一行 \texttt{print(href[i])}。下标从 0、编号从 1，故 \texttt{str(i + 1)}；数字与字符串不能直接 \texttt{+}：\texttt{'第' + 3} 报 \texttt{TypeError: can only concatenate str (not "int") to str}。\end{sloppypar}

\subsection{离线演练：样例源码的中间结果与输出}
两条样例（第 2 条来源前带 \texttt{<img>}、日期为相对时间）：清洗前 \texttt{title} 元素形如 \texttt{'\textbackslash n<em>阿里巴巴</em>马云:…\textbackslash n'}，\texttt{info} 形如 \texttt{'<img src="x.png">新浪财经\&nbsp;\&nbsp;3小时前'}；输出：
\texttt{1.电商趋势不可挡 水泥企业纷纷投入阿里巴巴、京东怀抱!(2018年12月13日 09:41-生意社)}，下一行网址；\texttt{2.阿里巴巴马云:创新是解决问题的唯一途径(3小时前-新浪财经)}，下一行网址。p\_title 不加 re.S 时 \texttt{title = []}，第五步什么也不打印——最常见的“\key{不报错但没结果}”。

\subsection{封装函数 baidu(company)：批量爬多家公司}
把第一至五步放进 \texttt{def baidu(company):}，网址末尾 \texttt{'...\&word=' + company}；定义 \texttt{companys = ['阿里巴巴', '万科集团', '百度', '腾讯', '京东']}；\texttt{for i in companys:} 中 \texttt{try: baidu(i); print(i + '百度新闻爬取成功')}，\texttt{except: print(i + '百度新闻爬取失败')}。\key{try/except}：try 中出错不崩溃，跳到 except 执行后继续下一家公司（异常处理详见第 7 章）。

\subsection{翻页规律 pn = (页码 − 1) × 10}\label{sec:D06-pn}
三页网址对比：第一页无 \texttt{pn}，第二页 \texttt{\&pn=10}，第三页 \texttt{\&pn=20}，只有末尾 \texttt{pn} 在变。每页 10 条，\texttt{pn} 是“跳过前多少条”，故 \(\texttt{pn}=(\text{页码}-1)\times 10\)（第一页相当于 \texttt{pn=0}）。单公司多页：\texttt{def baidu(page): num = (page - 1) * 10}，网址末尾 \texttt{+ '\&pn=' + str(num)}；\texttt{for i in range(10): baidu(i + 1)}（$i=0..9$ → 第 1..10 页），并打印 \texttt{'第' + str(i + 1) + '页爬取成功'}。

\end{multicols}

\subsection{随堂小测 1：多家公司 × 多页（def baidu(company, page)）}\label{sec:D06-quiz1}
题目：批量爬取四大行多页新闻，提示“用函数来定义 page、company”；交代码 + 运行视频（24 号晚 12 点前）。写法即 §\ref{sec:D06-full}：两参数函数，外层公司、内层页码（\texttt{range(1, 6)} 为第 1–5 页，同一公司各页连续），输出 20 行“××银行第 n 页爬取成功”并生成 txt。\textbf{易错}：\texttt{pn = page*10} 漏第 1 页（代入 page=1 应得 0）。\tagext\ 缺来源段时 info 比 title 短，\texttt{date[i]} 报 IndexError，先 \texttt{print(len(title), len(href), len(info))} 查等长。

\subsection{把结果写入文件：txt、csv 与 Excel \texorpdfstring{\tagext}{（拓展）}}
txt：\texttt{with open(文件名, 'a', encoding='utf-8') as f: f.write(字符串)}；\key{\texttt{'w'} 覆盖写、\texttt{'a'} 追加写}（多公司多页须 \texttt{'a'}）；write 不自动换行，需加 \texttt{'\textbackslash n'}；with 结束自动关闭。csv／Excel：
\begin{lstlisting}[style=pyD06]
import pandas as pd       # 多家公司：各次 df 收进列表，pd.concat(列表) 合并后再存
df = pd.DataFrame({'公司': company, '标题': title, '日期': date, '来源': source, '网址': href})  # 各列表长度须相同
df.to_csv('baidu_news.csv', index=False, encoding='utf-8-sig')  # utf-8-sig：Excel 打开不乱码
df.to_excel('baidu_news.xlsx', index=False)                      # 需安装 openpyxl
\end{lstlisting}

\subsection{实战提醒：页面改版、规则失效与“列表为空”排查顺序}
\texttt{c-title}、\texttt{c-author} 来自 2018 年前后的百度页面，改版后规则会失效——\key{方法不变，规则以当时 F12 源码为准}。“不报错但列表为空”：① \texttt{print(len(res))}，res 存文件后 Ctrl+F 搜公司名，源码无新闻 → 被拦截（查 UA、“百度安全验证”页）或动态加载（Selenium）；② 源码有新闻但 findall 为空 → 逐字核对文本 A/B（空格、引号、class 名）、是否漏 re.S；③ 列表不等长 → 部分新闻缺字段：改规则或先切大框再逐条提取。

% ======================================================================
\section{本章小结}
\begin{summary}[核心代码清单]
取源码 \texttt{requests.get(url, headers=headers).text}（UA 字典）→ 提取 \texttt{re.findall(p, res, re.S)} → “文本A\texttt{(.*?)}文本B”（\texttt{.*?} 填充、\texttt{(.*?)} 提取、须有右边界）→ 清洗 \texttt{strip()}、\texttt{replace}、\texttt{re.sub('<.*?>', '', s)}、\texttt{re.sub('[*]', '', s)}、\texttt{split('\&nbsp;\&nbsp;')} → 翻页 \texttt{pn = (page - 1) * 10} → \texttt{baidu(company, page)} + 双层循环 + try/except。
\end{summary}
\begin{english}[Key expressions]
\footnotesize
\begin{itemize}
\item \texttt{requests.get(url, headers=headers).text} returns the HTML string; without a User-Agent: empty page or 403.
\item \texttt{re.findall} returns a list of all non-overlapping matches (\texttt{[]} if none), only the group content if the pattern has a capturing group. \texttt{A(.*?)B} captures the text between boundaries A and B; an ungrouped \texttt{.*?} skips irrelevant content.
\item \texttt{re.S} (DOTALL) lets the dot match newlines; \texttt{re.sub('<.*?>', '', s)} removes all tags, whereas \texttt{str.replace} replaces a fixed substring; \texttt{'[*]'} matches a literal asterisk. Pagination: \texttt{pn = (page - 1) * 10}.
\end{itemize}
\end{english}

\section{易错汇总与数据采集流程设计模板}

\subsection{爬虫代码高频错误清单（找错）}\label{sec:D06-bugs}
\begin{center}\footnotesize\renewcommand{\arraystretch}{1.0}
\begin{tabularx}{\linewidth}{@{}p{7.6cm}L@{}}
\toprule
错误写法 & 正确写法／原因 \\
\midrule
\texttt{\{'User\_Agent': ...\}} & \texttt{'User-Agent'}（连字符）\\
\texttt{'\&pn=' + page * 10} & \texttt{str((page - 1) * 10)}：整数不能拼字符串，且少减 1 \\
\texttt{for page in range(3): baidu(company, page)} & 得 0,1,2：用 \texttt{page + 1} 或 \texttt{range(1, 4)} \\
\texttt{res = requests.get(url, headers=headers)} & 缺 \texttt{.text}（响应对象不是字符串）\\
\texttt{'<h3 class="c-title">.*?>(.*)</a>'}，无 re.S & 非贪婪 \texttt{(.*?)} 并加 \texttt{re.S} \\
\texttt{"日期：(.*?)"} & 无右边界得空串：\texttt{r"日期：(\textbackslash d\{4\}-\textbackslash d\{2\}-\textbackslash d\{2\})"} 或补右边界 \\
\texttt{re.sub('<.*>', '', t)} & \texttt{'<.*?>'}（贪婪会删掉标签间正文），再 \texttt{.strip()} \\
公司名 \texttt{'*ST国华'} 直接拼网址 & 先 \texttt{re.sub('[*]', '', company)} \\
无 try/except；结果不输出不保存 & 加 try/except 跳过失败页；print 或写文件 \\
\bottomrule
\end{tabularx}
\end{center}

\subsection{研究设计模板：用百度新闻构造公司—年度媒体关注度}\label{sec:D06-design}
例：媒体关注度能否降低股价崩盘风险（A 股，2015–2024）。\key{① 名单与网址}：A 股代码与简称（CSMAR/Wind），\texttt{re.sub('[*]', '', 简称)} 去星号，必要时加全称防同名；\texttt{word=简称\&pn=(page-1)*10}，\texttt{baidu(company, page)} 双层循环直到某页为空。\key{② 提取与清洗}：“文本A\texttt{(.*?)}文本B”+\texttt{re.S} 取四字段；去标签、split、strip；“X天前”“昨天”“YYYY年MM月DD日”统一为 YYYY-MM-DD；保留样本期；按网址去重。\key{③ 去噪与反爬}：逐条请求正文（timeout + try/except），候选容器或最长 div 抽正文，不含公司名者剔除；User-Agent + \texttt{time.sleep}，必要时 IP 代理，记录失败页事后补爬，动态加载/验证页改用 Selenium（第 7 章）。\key{④ 变量构造}：按公司—年度统计新闻条数 $N_{it}$，媒体关注度 $=\ln(1+N_{it})$；可按权威媒体来源构造稳健性指标，或做情感分析（第 8 章）得负面报道比例；与崩盘风险指标（NCSKEW、DUVOL）按代码—年度合并回归。
\textit{Pipeline: firm list → search URLs (word, pn) → GET + non-greedy regex → clean dates, drop noise → count by firm-year, $\log(1+N)$.}
